% ID: SEL-D0071-OA-E02
% ESTADO: cotejada. Segunda pasada visual y borrador compilado revisados.
% Fuente: Selectividad/sel_2009_matematicas_aplicadas/mat apli CC SS05.pdf
% SHA-256: ca51ab435c7584a8e047e19d0826f3273f3e88aed6110ccd00e46a40e85518b1
% Página: 1; opción A; ejercicio original 2.

Sea la función $f(x)=x^3-1$.
\begin{enumerate}[label=\alph*)]
\item \textbf{(1 punto)} Calcule los puntos de corte de la gráfica con los ejes, su monotonía y extremos relativos, si los tuviese.
\item \textbf{(1 punto)} Determine su curvatura y punto de inflexión.
\item \textbf{(1 punto)} Halle los puntos de la gráfica en los que la recta tangente tiene pendiente 3.
\end{enumerate}
